Nowadays, \gls{lidar} sensors are widely used in various applications, including autonomous vehicles, robotics, and environmental mapping [insert citation for robotics BAP and Pinliang book].
While the former two applications often require high-end and expensive \gls{lidar} systems, environmental mapping can be effectively performed using more affordable sensors.
However, even the cheapest commercial \gls{lidar} sensors, which were found during the research for this miniproject, cost approximately \$250[insert Unitree citation].
The primary motivation behind this mini-project (and the subsequent bachelor's thesis) is to develop a cost-effective platform for a $360\,\degree$ \gls{lidar} sensor that can be utilised for 3D spatial scanning and visualisation of the surrounding environment.
A subsequent bachelor's thesis will then focus on the \gls{dsp} techniques for the effective processing of the acquired point cloud data and proper utilisation of open-source visualisation tools (such as Python's Open3D library[insert citation of a site or a thesis for that library]) to recreate a 3D scene.